% Préambule LaTeX du guide timeseries, appliqué par pandoc (--include-in-header).
% Le Markdown reste la source ; ce fichier ne règle que la typographie.

\usepackage{microtype}          % césures et espacements fins
\usepackage{etoolbox}
\usepackage{booktabs}
\usepackage{array}
\usepackage{fancyhdr}
\usepackage{titlesec}

% --- Titres : alignés à gauche, sans exception ---------------------------
\titleformat{\section}{\Large\bfseries\raggedright}{\thesection}{1em}{}
\titleformat{\subsection}{\large\bfseries\raggedright}{\thesubsection}{1em}{}
\titleformat{\subsubsection}{\bfseries\raggedright}{\thesubsubsection}{1em}{}
\setlength{\parindent}{0pt}
\setlength{\parskip}{0.6em}

% --- Les tableaux comparatifs sont larges : on les resserre --------------
\AtBeginEnvironment{longtable}{\small}
\AtBeginEnvironment{tabular}{\small}

% --- En-tête et pied de page --------------------------------------------
\pagestyle{fancy}
\fancyhf{}
\fancyhead[L]{\small\nouppercase{\leftmark}}
\fancyhead[R]{\small timeseries}
\fancyfoot[C]{\small\thepage}
\renewcommand{\headrulewidth}{0.4pt}

% --- Exposants : Latin Modern ne les a pas, on les fabrique -------------
\usepackage{newunicodechar}
\newunicodechar{⁵}{\textsuperscript{5}}
\newunicodechar{²}{\textsuperscript{2}}
\newunicodechar{ᵉ}{\textsuperscript{e}}
\newunicodechar{ʳ}{\textsuperscript{r}}
\newunicodechar{·}{\textperiodcentered}
\newunicodechar{±}{$\pm$}
\newunicodechar{×}{$\times$}
\newunicodechar{−}{$-$}

% --- Les blocs de code gardent leurs caractères semi-graphiques ----------
\usepackage{fvextra}
\fvset{breaklines=true, breakanywhere=true, fontsize=\small}
